
\chapter{Цифровой рубль}\label{ch:cbdc}
\highlight{Цифровой рубль~--- третья форма денег наряду с~наличными
и~безналичными} (см.~раздел~\ref{sec:banks-forms-of-money}):
цифровая форма национальной валюты, которую Банк России эмитирует
напрямую~\cite{fz-340,cbr-digital-ruble}, с~отдельным реестром
обязательств на~платформе Банка России. Цифровой рубль относится
к~классу CBDC (\textit{central bank digital currency}, цифровая
валюта центрального банка) и~представляет собой розничную CBDC:
платформа рассчитана на~переводы между гражданами, ТСП и~бюджетными
распорядителями, в~отличие от~оптовых (\textit{wholesale}) CBDC,
рассчитанных на~расчёты между финансовыми
организациями~\cite{bis-wp880-cbdc-architectures}. Сравнение
с~СБП~--- в~разделе~\ref{sec:sbp-cbdc}.

\section{Двухуровневая архитектура}\label{sec:cbdc-architecture}

Банк России эмитирует цифровой рубль и~управляет платформой;
участники платформы (банки и~Федеральное казначейство) открывают
клиентам кошельки и~исполняют поручения на~платформе Банка России.
Граждане и~ТСП работают с~кошельком через мобильное приложение
банка; главные распорядители бюджетных
средств (ГРБС)\footnote{ГРБС~--- орган, определённый ст.~158
Бюджетного кодекса~РФ~\cite{bk-rf}, распределяющий бюджетные
ассигнования и~лимиты обязательств между подведомственными
получателями.} работают через интерфейсы
Казначейства~\cite{cbr-digital-ruble}.
Требования к~защите информации для~участников платформы установлены
Положением Банка России от~07.12.2023~№~833-П \enquote{О~защите
информации для~участников платформы цифрового рубля}~\cite{cbr-833p}.

\begin{figure}[tbp]
  \centering
  \includefiguresvg[width=\linewidth]{assets/figures/ch23-cbdc/two-tier-model}
  \caption{Двухуровневая архитектура платформы цифрового
  рубля; сценарий выплаты от~ГРБС физическому лицу.}\label{fig:cbdc-two-tier}
\end{figure}
\vspace{-4pt}

Платформа Банка России включает (рис.~\ref{fig:cbdc-two-tier})
эмиссию цифрового рубля, единый реестр кошельков и~операций,
централизованный расчётный движок и~узлы DLT (\textit{distributed
ledger technology}, технология распределённого реестра; формальное
определение~--- в~разделе~\ref{sec:cbdc-blockchain-fork}), реплицирующие
журнал операций между утверждёнными участниками. На~втором уровне
работают участники платформы: банки обслуживают розничных
и~корпоративных клиентов, Федеральное казначейство ведёт бюджетные
кошельки и~целевые выплаты с~казначейским сопровождением; круг
может расширяться по~решению Банка России. Всю предварительную
работу до~поступления распоряжения на~платформу участник выполняет
в~собственном лицензионном периметре: идентификация
клиента, проверки по~ПОД/ФТ (противодействие отмыванию доходов
и~финансированию терроризма) и~антифроду, проверка лимитов
и~реквизитов. На~платформу поступает уже подписанное
и~проверенное распоряжение, готовое к~расчёту.

Отличие от~безналичного рубля лежит в~природе обязательства.
Безналичный рубль~--- обязательство коммерческого банка перед
клиентом; цифровой рубль~--- прямое обязательство Банка России.
Для~клиента исчезает кредитный риск банка; для~банков возникает
вопрос о~судьбе депозитной базы~\cite{cbr-digital-ruble}.

\section{Выбор между приватным и~публичным блокчейном}\label{sec:cbdc-blockchain-fork}

\emph{Распределённый реестр} (\textit{distributed ledger},
DLT)~--- база данных, копии которой одновременно хранятся
у~нескольких независимых участников. Когда нужно записать новую
операцию, участники договариваются между собой, что запись корректна,
и~каждый обновляет свою копию. \emph{Блокчейн}~--- частный вид
распределённого реестра, где записи объединены в~цепочку блоков,
а~каждый блок ссылается на~предыдущий через криптографический хеш;
изменить старую запись задним числом нельзя без переподписания всей
цепочки. \emph{Валидатор}~--- участник с~правом подтверждать
новые записи; механизм, которым валидаторы приходят к~согласию,
называется \emph{консенсусом}.

\emph{Proof-of-work} (PoW, доказательство выполненной работы)
заставляет валидаторов соревноваться в~подборе хеша, удовлетворяющего
условию; победитель тратит реальное электричество и~машинное время,
и~именно эти затраты делают подделку невыгодной, как это устроено
в~Bitcoin.
\emph{Proof-of-stake} (PoS, доказательство доли владения) отдаёт
право подписать блок участнику, заблокировавшему в~залог крупную
сумму криптовалюты; за~некорректное поведение залог конфискуется
(механизм \textit{slashing}), и~по~такой схеме с~2022~года работает
Ethereum.
\emph{Proof-of-authority} (PoA, доказательство полномочия) допускает
к~подписи блоков только тех участников, которым оператор сети заранее
выдал такие полномочия; ни~работы, ни~залога не~требуется,
доверие держится на~репутации и~договоре. PoA~--- частый выбор
для~приватных корпоративных и~банковских сетей, включая платформы CBDC.

\emph{Публичный} (\textit{permissionless}) реестр~--- такой, где
валидатором может стать кто угодно, набрав ресурс по~правилам
PoW или~PoS. \emph{Приватный} (\textit{permissioned}) реестр~---
такой, где список валидаторов утверждён заранее (PoA или~его
варианты), и~войти в~него можно только по~разрешению оператора
сети. На~публичной сети любой пользователь видит все транзакции;
на~приватной видимость определяется правилами оператора.

В~консультативной фазе Банк России сравнивал три варианта:
централизованную базу данных под~полным контролем регулятора,
децентрализованную сеть на~базе распределённого реестра
и~гибридную архитектуру~\cite{cbr-concept-2021}. В~выбранной
гибридной модели распределённый реестр работает между узлами
Банка России и~банков-участников, но~все эти узлы принадлежат
проверенному кругу, а~окончательный расчёт по~операциям выполняют
централизованные компоненты Банка России~\cite{cbr-concept-2021}.
Полностью централизованный вариант проигрывал в~устойчивости (одна
точка отказа), полностью децентрализованный~---
в~производительности: DLT медленнее обычной СУБД.

Между приватным и~публичным реестром регуляторы выбирают
единогласно. \highlight{Ни~один центральный банк, экспериментирующий
с~CBDC, не~выбрал публичный вариант}: BIS (\textit{Bank for
International Settlements}, Банк международных расчётов в~Базеле),
проанализировав все доступные на~2020~год отчёты регуляторов,
не~нашёл ни~одного~\cite{bis-wp880-cbdc-architectures}.

Юридическая финальность платежа. В~PoW транзакция считается
\enquote{достаточно надёжной} лишь после того, как поверх неё
накопится несколько блоков; до~этого её теоретически можно отменить,
перестроив цепочку. PoS снимает остроту проблемы, но~полностью её
не~устраняет: при~разделении сети возможны откаты. Платёжное право
требует, чтобы в~момент списания деньги были переведены
\emph{окончательно}, без~оговорок про~вероятность отката. PoA
с~доверенным списком валидаторов даёт такую финальность
по~построению: подпись утверждённого узла сразу фиксирует запись.

Пропускная способность. Bitcoin (PoW) обрабатывает порядка семи
операций в~секунду на~всю сеть, Ethereum после перехода на~PoS
на~уровне основной цепи (\textit{mainnet})~--- порядка пятнадцати;
розничная платёжная инфраструктура страны должна выдерживать тысячи
в~секунду круглосуточно. Открытый консенсус так устроен: чем больше
валидаторов и~чем меньше у~оператора рычагов их отбирать, тем дороже
договариваться о~каждой записи. PoA с~небольшим утверждённым кругом
участников снимает это ограничение, потому что согласование
сводится к~подписям заранее известных сторон.

Эмиссия. Цифровой рубль~--- обязательство Банка России, и~право
выпустить новый рубль есть у~Банка России. На~публичной сети любой
валидатор, собравший достаточный ресурс по~правилам PoW или~PoS,
в~принципе может влиять на~правила выпуска, что для~национальной
валюты неприемлемо.

\highlight{Поэтому в~архитектуре цифрового рубля DLT остаётся
способом надёжно реплицировать регистр между утверждёнными узлами
и~зафиксировать историю операций криптографически}. Слово
\enquote{блокчейн} в~этом контексте означает закрытую сеть
с~подписями участников из~заранее известного списка. Та же модель
применяется в~китайском e-CNY (\textit{Digital Currency Electronic
Payment}, DC/EP~--- розничная CBDC Народного банка Китая
с~двухуровневой архитектурой и~закрытым кругом
валидаторов)~\cite{pboc-ecny-2021}. Технологически разница такого
реестра с~хорошо защищённой банковской базой данных невелика:
отличаются только репликация и~формат подписи, открытости участия
нет ни~там, ни~здесь.

Окончательный расчёт по~операциям выполняют не~валидирующие узлы
DLT, а~централизованные компоненты процессинга Банка
России~\cite{cbr-concept-2021}. Bitcoin ограничен семью операциями
в~секунду из-за стоимости согласования между валидаторами открытого
консенсуса. В~гибридной архитектуре цифрового рубля этой стоимости
нет: согласование подменено подписью утверждённого узла, а~сам
расчёт идёт по~логике обычной СУБД с~горизонтальным
масштабированием. Распределённый реестр на~этом фоне выполняет
роль аудиторского журнала; расчётным движком он не~служит. Целевые
TPS (\textit{transactions per second}, операций в~секунду) Банк
России не~публикует.

ПОД/ФТ, антифрод, проверку лимитов и~реквизитов, идентификацию
клиента выполняют участники платформы в~своём
лицензионном периметре до~того, как распоряжение доходит
до~Банка России (см.~раздел~\ref{sec:cbdc-architecture}). Ядро видит
уже подписанное и~проверенное распоряжение~--- типовую операцию
\enquote{списать с~одного кошелька, зачислить на~другой}
в~транзакции с~гарантиями ACID (\textit{atomicity, consistency,
isolation, durability}~--- атомарность, согласованность,
изолированность, устойчивость; классический набор свойств
транзакций реляционной СУБД). С~такой нагрузкой давно справляются
ядро НСПК и~инфраструктура СБП, обрабатывающие тысячи операций
в~секунду на~всю платёжную систему страны.

Пилот с~августа 2023~года и~до~конца 2024-го дал более 38~тыс.\
переводов между физическими лицами за~полтора
года~\cite{cbr-onrnps-2025}; с~сентября 2025~года в~Татарстане,
Чувашии и~Ростовской области
отрабатываются целевые выплаты от~ГРБС физическим лицам
(рис.~\ref{fig:cbdc-two-tier}). По~Федеральному закону
от~23.07.2025~№~248-ФЗ \enquote{О~поэтапном внедрении цифрового
рубля} обязанность принимать цифровой рубль наступает в~три
волны: с~1~сентября 2026~--- продавцы с~выручкой свыше
120~млн руб./год, с~1~сентября 2027~--- свыше 30~млн руб./год,
с~1~сентября 2028~--- все остальные, кроме точек с~выручкой
до~5~млн руб./год и~точек без~интернета~\cite{fz-248}. Розничный
безналичный поток в~III~квартале 2025~года составил около
2,8~тыс.\ операций в~секунду на~всю страну~\cite{cbr-payments-2025q3};
до~этого уровня платформе цифрового рубля остаются три года
на~нагрузочное тестирование и~рост инфраструктуры.

\section{Смарт-контракты}\label{sec:cbdc-smart-contracts}

\emph{Смарт-контракт} в~контексте цифрового рубля~--- программный
сценарий, опубликованный на~платформе Банка России и~автоматически
исполняющий перевод цифровых рублей при~наступлении заданных условий
(дата, событие, подтверждение от~стороны сделки) без~вмешательства
оператора~\cite{cbr-digital-ruble}. Возможность встроить условия
исполнения непосредственно в~единицу денег называют
\emph{программируемыми деньгами} (\textit{programmable money});
у~безналичного рубля такого свойства нет~\cite{bis-wp880-cbdc-architectures}.

Простые смарт-контракты уже работают в~пилоте
и~автоматически переводят цифровые рубли в~заданные дату и~время.
К~маю 2025~года в~пилоте исполнено более 17~тыс.\ таких
контрактов~\cite{cbr-digital-ruble}; типовые сценарии~--- автоплатёж
за~жилищно-коммунальные услуги, автопополнение кошелька, поэтапная
оплата подрядчику~\cite{gazprombank-smart-contracts-2025}.

Коммерческая платформа смарт-контрактов, на~которой банки смогут
публиковать собственные сценарии с~конфиденциальными условиями
и~юридической обёрткой, находится в~стадии разработки. В~марте
2026~года глава Банка России сообщила, что регулятор
\enquote{представит для~обсуждения концепцию платформы коммерческих
смарт-контрактов} в~течение
2026~года~\cite{cbr-nabiullina-abr-2026}. Целевые выплаты,
страховые сценарии и~B2B-расчёты с~условиями пока остаются
на~стадии концепции.

\section{Цифровой рубль и~СБП}\label{sec:cbdc-vs-sbp}

Цифровой рубль и~СБП часто сравнивают: оба строит государство, оба
обещают более дешёвый платёж. Архитектурно это разные системы: СБП
переводит безналичные рубли между банковскими счетами, а~цифровой
рубль хранится в~отдельном кошельке как самостоятельная форма
денег~\cite{cbr-digital-ruble}.

\vspace{-6pt}
\begin{table}[H]
  \caption{СБП и~цифровой рубль на~одном уровне сравнения.}\label{tab:cbdc-vs-sbp}
  \begingroup
  \footnotesize
  \setlength{\tabcolsep}{6pt}
  \renewcommand{\arraystretch}{1.2}
  \begin{tabularx}{\textwidth}{
      @{}P{0.24\textwidth} Y Y@{}
    }
    \toprule
    \headrow
    \thd{} & \thd{СБП} & \thd{Цифровой рубль} \\
    \midrule
    Природа обязательства & Коммерческого банка & Банка России \\
    Оператор реестра & Банк-участник & Банк России \\
    Расчётный цикл & Круглосуточно, в~реальном времени & Круглосуточно, в~реальном времени \\
    Тариф для~ТСП (С2B) & до~0,7\,\% (по~категориям) & 0\,\% в~пилоте; льготный
      тариф на~старте обязательного приёма \\
    Офлайн-платёж & Нет & Предусмотрен архитектурой, в~пилоте не~запущен \\
    Смарт-контракты & Нет & Простые~--- в~пилоте; коммерческая
      платформа в~стадии разработки \\
    \bottomrule
  \end{tabularx}
  \endgroup
\end{table}

Для~клиента и~ТСП оба сценария выглядят одинаково: быстрый перевод
по~номеру телефона или~QR-коду. Различия проявляются в~распределении
операционного и~кредитного риска. В~СБП клиент несёт кредитный риск
банка-участника, в~цифровом рубле этот риск исчезает по~построению.
Для~банка-участника СБП остаётся источником комиссии и~бренда,
а~цифровой рубль~--- источником оттока депозитной базы, который
оператор должен компенсировать другими доходами.

\section{Цифровой рубль в~международном контексте}\label{sec:cbdc-international}

\highlight{Россия в~2026~году~--- единственная страна, перешедшая
от~пилота к~законодательно обязательному приёму CBDC на~рознице}~\cite{atlantic-council-cbdc-tracker}.

\needspace{4\baselineskip}
\subsection*{e-{CNY} (КНР)}

К~ноябрю 2025~года в~Китае накоплено~3,48~млрд транзакций в~e-CNY
на~общую сумму~16,7~трлн юаней (около~\$2,37\,трлн)~--- рост
от~2023~года более чем в~восемь раз~\cite{gov-cn-ecny-2026}. Пилот
действует в~29~регионах. С~1~января 2026~года Народный банк Китая ввёл
обновлённую модель управления, в~которой e-CNY переходит от~аналога
наличных к~\enquote{digital deposit money} с~депозитными функциями;
это приближает его к~стандартной банковской инфраструктуре. e-CNY
остаётся розничной CBDC с~закрытым кругом валидаторов, без~публичного
блокчейна (см.~раздел~\ref{sec:cbdc-blockchain-fork}).

\needspace{4\baselineskip}
\subsection*{Digital euro (ЕС)}

Подготовительная фаза ЕЦБ шла с~ноября 2023 по~октябрь 2025~года;
30~октября 2025~года Управляющий совет принял решение о~переходе
в~следующую фазу~\cite{ecb-digital-euro-next-phase-2025}.
Целевой запуск~--- 2029~год при~условии принятия европейскими
со-законодателями Регламента о~цифровом евро в~2026~году;
промежуточный пилот ожидается в~середине 2027~года. Rulebook~---
свод единых правил для~эмиссии, держателей, ТСП и~инфраструктурных
провайдеров~--- прошёл четырёхмесячную market consultation
и~зафиксирован в~итоговом отчёте подготовительной
фазы~\cite{ecb-deprp-2025}.

\needspace{4\baselineskip}
\subsection*{Digital pound (Великобритания)}

Банк Англии прошёл консультацию 2023~года и~работает в~design phase
до~2026~года. В~августе 2025~года открыт Digital Pound Lab~---
экспериментальная платформа для~отрасли, где тестируются
сценарии и~бизнес-модели~\cite{boe-digital-pound-design-2025}.
Решение запускать или~отказываться от~проекта Банк Англии
вместе с~HM~Treasury планирует представить в~2026~году; запуск
требует отдельного парламентского решения. Великобритания, в~отличие
от~ЕС, не~вводит регулирование до~технологической готовности.

\needspace{4\baselineskip}
\subsection*{США}

Указ Президента США \enquote{Strengthening American Leadership
in Digital Financial Technology} от~23~января 2025~года
запретил федеральным агентствам выпускать CBDC, в~том числе
Федеральной резервной системе~\cite{trump-eo-digital-assets-2025}.
Это разворот политики предыдущей администрации, при~которой
ФРС разрабатывал Project Cedar. В~июле 2025~года Anti-CBDC
Surveillance State Act (H.R.~1919) Тома Эммера прошёл Палату
представителей; закон закрепляет запрет федеральным законом,
а~не~указом~\cite{anti-cbdc-act-2025}. Параллельно курс смещается
к~долларовым стейблкоинам как частной альтернативе CBDC.

Китай форсирует развёртывание с~прицелом на~глобальную инфраструктуру;
ЕС идёт через нормативное оформление, не~опережая политическое
решение; Великобритания удерживает право не~запускать проект вообще;
США отказались от~CBDC через прямой правовой запрет. Россия идёт
ортогонально: закон~248-ФЗ~\cite{fz-248} превращает приём в~обязанность
ТСП тремя волнами с~1~сентября 2026 по~1~сентября 2028~года, и~эта
обязательность~--- то, чего нет ни~у~одного из~перечисленных проектов.
Международной совместимости цифрового рубля с~e-CNY, евро или~фунтом
ждать не~приходится: платформы строятся под~разные правила,
и~синхронизация требует отдельного двустороннего соглашения.

\section{{Project} {mBridge} и~Project {Agor{\'a}}}\label{sec:cbdc-mbridge-agora}

\highlight{Инфраструктура трансграничных CBDC раскалывается
на~две непересекающиеся группы центральных банков}. Project
mBridge~--- мультивалютная оптовая платформа на~базе DLT
для~прямых расчётов между центральными банками~--- запущена
в~2021~году как совместный проект PBoC, HKMA, Bank of Thailand,
Central Bank of the UAE и~BIS Innovation Hub. В~2024~году
к~проекту присоединился Saudi Central Bank, и~mBridge вышел
из~стадии proof-of-concept в~minimum viable product.

31~октября 2024~года генеральный директор BIS Агустин Карстенс
объявил о~выходе BIS из~проекта с~формулировкой \enquote{graduation}~---
проект достиг зрелости, при~которой партнёрам участие BIS более
не~требуется~\cite{bis-mbridge-graduation-2024}. Это произошло
через две недели после саммита БРИКС в~Казани, где mBridge упоминался
как инструмент де-долларизации. BIS подчеркнул, что mBridge
\enquote{не~является BRICS-bridge} и~что переход к~независимой
работе не~связан с~политическим давлением. Партнёры~--- PBoC, HKMA,
центробанки Таиланда, ОАЭ и~Саудовской Аравии~--- продолжают
платформу самостоятельно.

Параллельно BIS усиливает Project Agorá. В~составе~--- Bank of
France, Bank of Japan, Bank of Korea, Bank of Mexico, Swiss
National Bank, Bank of England и~Federal Reserve Bank of New
York~\cite{bis-project-agora}. Цель~--- унифицированный леджер
для~tokenised commercial bank money и~wholesale central bank money
с~сохранением контролей корреспондентского банкинга
(\textit{correspondent banking}). В~Agorá нет участников из~БРИКС.

Российский трансграничный набор инструментов~--- СПФС, БРИКС~Pay,
потенциально mBridge через партнёров в~Азии~--- к~Agorá подключения
не~получает. Двусторонние интеграции цифрового рубля с~e-CNY
или~другими CBDC участников mBridge технически возможны; интеграция
с~CBDC участников Agorá упирается в~санкции и~требования
корреспондентского банкинга, которые Agorá сохраняет по~построению.

\section{{DCash} в~Восточно-Карибском валютном союзе}\label{sec:cbdc-dcash}

12~января 2024~года Центральный банк Восточно-Карибского валютного
союза (ECCB) закрыл пилот DCash после~34~месяцев работы~--- редкий
публичный случай отзыва розничной
CBDC~\cite{eccb-dcash-shutdown-2024}. В~октябре 2025~года ECCB
запустил тендер DCash~2.0 с~расширенным набором сценариев:
P2P, P2B, B2B, G2P, P2G~\cite{eccb-dcash-2-rfvi-2025}. На~112-м
заседании Монетарного совета ECCB 13~февраля 2026~года разработка
DCash~2.0 приостановлена в~пользу межбанковской Fast Payment System
и~участия в~региональном пилоте CARICOM Payments and Settlement
System~\cite{eccb-mc-112-2026}.

Пилот работал без~обязательного приёма; держателю DCash было удобнее
использовать привычные карточные и~электронные кошельки.
Без~массового сценария платформа не~накапливала транзакционный
поток: за~34~месяца совокупный объём остался на~уровне пилотного,
несопоставимом со~штатной розничной нагрузкой региона. У~российского
пилота цифрового рубля похожая статистика на~ту же стадию:
38~тыс.\ переводов P2P за~полтора года~\cite{cbr-onrnps-2025}.
Разница в~том, что у~российского цифрового рубля впереди закон
248-ФЗ с~обязательным приёмом, у~DCash такого рычага не~было.

\highlight{Розничная CBDC без~обязательного приёма не~накапливает
критическую массу}. e-CNY масштабируется через государственные
сценарии (зарплаты, целевые выплаты, госуслуги) и~интеграцию
с~Alipay и~WeChat~Pay; цифровой рубль~--- через законодательное
обязательство приёма. Пилоты без~якорного сценария повторяют
траекторию DCash.
